% =============================================================================
% mc_tables.tex — compiled view of all xtife Monte Carlo tables (JSS draft)
% Build:  pdflatex mc_tables.tex   (run twice for page refs)
% Source .tex tables are generated by collate_mc.R; do not edit them by hand.
% =============================================================================
\documentclass[10pt]{article}
\usepackage[a4paper,landscape,margin=1.4cm]{geometry}
\usepackage{amsmath}
\usepackage{amssymb}
\usepackage{booktabs}
\usepackage{caption}
\captionsetup{font=small,labelfont=bf,width=.95\linewidth}
\usepackage{etoolbox}
% shrink all tabulars slightly so the wide tables fit the landscape page
\AtBeginEnvironment{tabular}{\footnotesize\setlength{\tabcolsep}{4pt}}
\setcounter{topnumber}{2}
\renewcommand{\topfraction}{.95}
\renewcommand{\textfraction}{.02}
\renewcommand{\floatpagefraction}{.9}

\title{\textbf{xtife: Monte Carlo Simulation Results}\\[2mm]
\large Interactive Fixed Effects Estimators --- Bai (2009), Moon \& Weidner (2017), SWW (2025)}
\author{B = 1000 replications per cell; $r_{\mathrm{true}} = 2$;
$\beta_{\mathrm{true}} = 1$ (static) / $0.5$ (dynamic)}
\date{\today}

\begin{document}
\maketitle

\subsection*{Data generating processes}

We report simulation results based on $B=1000$ replications of the panel model
\[
y_{it} \;=\; \beta\, x_{it} \;+\; \sum_{\rho=1}^{2}\lambda_{i\rho}\,f_{t\rho} \;+\; u_{it},
\qquad i=1,\dots,N,\; t=1,\dots,T,
\]
with $r_{\mathrm{true}}=2$ interactive factors and no additive fixed effects
($\alpha_i=\xi_t=0$). Throughout, the loadings are
$\lambda_{i\rho}\stackrel{iid}{\sim} N(0,1)$, independent across $i$ and
$\rho$ and of all other primitives. The following four DGPs are employed.

\paragraph{DGP 1: static balanced panel.}
We set $\beta_0 = 1$ and, for $\rho\in\{1,2\}$,
$f_{t\rho}\stackrel{iid}{\sim}N(0,1)$ and
$\mu_{i\rho}\stackrel{iid}{\sim}N(0,1)$, with $\mu_i$ independent of
$\lambda_i$. The regressor is
\[
x_{it} \;=\; \sum_{\rho=1}^{2}(\lambda_{i\rho}+\mu_{i\rho})\,f_{t\rho}
\;+\; \varepsilon_{it},
\qquad \varepsilon_{it}\stackrel{iid}{\sim}N(0,1),
\]
so $x_{it}$ is strictly exogenous with respect to $u_{it}$ but loads on the
same factors as $y_{it}$: pooled OLS that ignores the factor structure is
biased (the OLS column, $\approx 0.4$). The independent extra loading
$\mu_i$ ensures $x$ is not an exact linear combination of
$(f_{t1},f_{t2})$; otherwise the projection of $y$ on $x$ would absorb one
factor direction and the information criterion would select $\hat r = 1$.
Errors $u_{it}$ follow the homoskedastic or heteroskedastic specification
described below. The panel is fully observed.

\paragraph{DGP 2: static unbalanced panel.}
Identical to DGP 1, after which cells are removed completely at random:
a subset of $\lfloor \pi N T\rfloor$ index pairs $(i,t)$ is drawn uniformly
without replacement and retained, $\pi\in\{0.8,\,0.6\}$ (the ``fill'' rate).
The missingness indicator $d_{it}$ is therefore independent of
$(\lambda_i, f_t, u_{is}, x_{is})$ for all $(i,s)$ --- SWW's missing
Pattern~1 with a homogeneous observation probability.

\paragraph{DGP 3: dynamic balanced panel.}
We set $\beta_0=0.5$ and let the factors follow stationary AR(1) processes,
\[
f_{t\rho} \;=\; 0.5\, f_{t-1,\rho} + \eta_{t\rho},
\qquad \eta_{t\rho}\stackrel{iid}{\sim} N\!\bigl(0,\,1-0.5^{2}\bigr),
\qquad f_{1\rho}\sim N(0,1),
\]
so that $\mathrm{Var}(f_{t\rho})=1$ for all $t$. The outcome is generated
recursively over $t=1,\dots,T_0+T$ with burn-in length $T_0=50$:
\[
y_{i1} = \lambda_i' f_1 + u_{i1},
\qquad
y_{it} = 0.5\, y_{i,t-1} + \lambda_i' f_t + u_{it}, \quad t\ge 2,
\]
and the regressor is the lagged outcome, $x_{it} = y_{i,t-1}$. The first
$T_0+1$ periods are discarded, leaving an estimation sample of $T-1$ time
periods (each retained $y_{it}$ paired with its observed lag). The errors
$u_{it}$ are \emph{serially uncorrelated} in both variants: with a lagged
dependent variable, serial correlation in $u_{it}$ would render
$y_{i,t-1}$ contemporaneously endogenous
($\mathbb{E}[x_{it}u_{it}]\neq 0$), which no incidental-parameter bias
correction can repair. Thus $x_{it}$ is predetermined (weakly exogenous):
$\mathbb{E}[u_{it}\,|\,y_{i,t-1},y_{i,t-2},\dots,\lambda_i,f]=0$.
This design generates the Nickell-type $O(1/T)$ incidental-parameter bias,
which halves as $T$ doubles (visible in the IFE column of Table~3).

\paragraph{DGP 4: dynamic unbalanced panel.}
Identical to DGP 3, after which cells are removed completely at random
exactly as in DGP 2 ($\lfloor \pi N (T-1)\rfloor$ cells retained,
$\pi\in\{0.8,0.6\}$). Missing lags are handled by the unbalanced estimator's
observed-cell objective; no imputation of $x$ is performed.

\paragraph{Error variants.}
In the \emph{homoskedastic} tables, $u_{it}\stackrel{iid}{\sim}N(0,1)$.
In the \emph{heteroskedastic} tables, $u_{it}=\sigma_{it}\,e_{it}$ with
$e_{it}\stackrel{iid}{\sim}N(0,1)$ and
\[
\sigma_{it} \;=\; c\,\Bigl(0.4 + 0.8\,
\frac{|\lambda_i' f_t|}{\mathrm{sd}(\lambda' f)}\Bigr),
\]
where $\mathrm{sd}(\lambda' f)$ is the in-sample standard deviation of the
factor component and the constant $c$ rescales $\sigma_{it}$ so that
$\frac{1}{NT}\sum_{it}\sigma_{it}^2 = 1$, keeping the overall noise level
identical to the homoskedastic case. Because $\sigma_{it}$ rises with
$|\lambda_i'f_t|$ --- the same component the regressor loads on --- the
homoskedastic standard SE undercovers by design, and valid inference
requires the robust/cluster/HAC variants.

\paragraph{Estimation.}
For every replication we compute:
(i) pooled OLS ignoring the factor structure (grand-mean demeaned;
for the unbalanced designs, OLS with intercept on the observed cells);
(ii) the raw IFE estimator with $r=2$ treated as known ---
\texttt{ife()} for DGPs 1 and 3 (the Bai 2009 least-squares estimator;
\texttt{method = "dynamic"} for DGP 3) and \texttt{ife\_unbalanced()} for
DGPs 2 and 4 (the SWW 2025 EM estimator with nuclear-norm-regularised
initialisation, \texttt{init = "nnr"}); and
(iii) the bias-corrected estimator:
Bai (2009) $B/N + C/T$ for DGP 1;
SWW $b_3$--$b_6$ under strict exogeneity for DGP 2;
Moon--Weidner (2017) $B_1,B_2,B_3$ with lag bandwidth $M_1=1$ for DGP 3;
and the SWW weak-exogeneity correction $b_2$--$b_6$ for DGP 4, where the
dynamic term $b_2$ and the HAC variance use a Bartlett kernel with bandwidth
$L_T=\lfloor 2\,T^{1/5}\rfloor$.
Confidence intervals for the raw estimator are reported with homoskedastic
standard, HC1-robust, cluster-robust (by unit), and --- for the unbalanced
estimator --- HAC standard errors; the bias-corrected estimator is paired
with the SE matching the error structure (standard under homoskedasticity,
HC1-robust under heteroskedasticity, HAC for DGP 4).
The number of factors is selected by a BIC-type information criterion over
$r\in\{0,1,\dots,5\}$; selection accuracy is reported as $P(\hat r=r)$ in
Tables~1--4 and in full in Table~5.
Replications use the L'Ecuyer--CMRG generator with a deterministic seed per
$(\mathrm{DGP},N,T,\pi,b)$ cell, so every entry is exactly reproducible.

\paragraph{Grid.}
$N\in\{30,50,100,200\}$, $T\in\{10,50,100,200\}$ ($T=10$ omitted for the
dynamic designs, which require the burn-in), and fill
$\pi\in\{80\%,60\%\}$ for DGPs 2 and 4. Note that the loadings are
mean-zero and the missingness non-informative, so the intercept-free
specification (\texttt{force = "none"}) is correctly specified throughout;
this deliberately differs from the simulation design of SWW (2025), which
uses mean-one loadings, two regressors, and factor-correlated missing
patterns.

\subsection*{Notation and definitions}

Each cell of Tables 1--4 summarises $B=1000$ Monte Carlo replications.
Let $\hat\beta_b$ denote the estimate in replication $b=1,\dots,B$,
$\beta_0$ the true coefficient,
$\bar{\hat\beta} = B^{-1}\sum_{b}\hat\beta_b$,
and $\widehat{\mathrm{se}}_b$ the estimated standard error in replication $b$.
The reported statistics are:

\begin{itemize}\setlength\itemsep{1pt}
  \item \textbf{Bias} (columns OLS, IFE, BC):
    $\mathrm{Bias} = B^{-1}\sum_{b=1}^{B}\bigl(\hat\beta_b-\beta_0\bigr)$,
    computed separately for pooled OLS (ignoring the factor structure),
    the raw IFE estimator with $r$ treated as known, and the analytically
    bias-corrected (BC) IFE estimator.
  \item \textbf{SD}: Monte Carlo standard deviation of the raw IFE estimator,
    $\mathrm{SD} = \bigl[(B-1)^{-1}\sum_{b}(\hat\beta_b-\bar{\hat\beta})^2\bigr]^{1/2}$.
  \item \textbf{RMSE / BC.RMSE}:
    $\mathrm{RMSE} = \bigl[B^{-1}\sum_{b}(\hat\beta_b-\beta_0)^2\bigr]^{1/2}
    = \sqrt{\mathrm{Bias}^2+\tfrac{B-1}{B}\mathrm{SD}^2}$,
    for the raw and the bias-corrected estimator respectively.
  \item \textbf{SE/SD}: calibration ratio
    $B^{-1}\sum_{b}\widehat{\mathrm{se}}_b \,/\, \mathrm{SD}$
    using the homoskedastic standard SE of the raw estimator.
    Values near 1 indicate that the analytical SE formula replicates the
    true sampling variability; values below 1 indicate underestimation
    (expected for the \emph{standard} SE under heteroskedastic errors).
  \item \textbf{95\% Coverage} (Std / Rob / Cl / HAC): empirical coverage
    of the nominal 95\% confidence interval for the \emph{raw} estimator,
    $B^{-1}\sum_{b}\mathbf{1}\bigl\{|\hat\beta_b-\beta_0|\le
    1.96\,\widehat{\mathrm{se}}_b\bigr\}$,
    where $\widehat{\mathrm{se}}_b$ is, respectively, the homoskedastic
    standard SE, the HC1 heteroskedasticity-robust SE, the cluster-robust SE
    (clustered by unit), or the Bartlett-kernel HAC SE with bandwidth
    $L_T=\lfloor 2T^{1/5}\rfloor$ (unbalanced estimator only).
  \item \textbf{BC.Std / BC.Rob / BC.HAC}: coverage of the same nominal 95\%
    interval centred at the \emph{bias-corrected} estimator, paired with the
    SE matching the error structure: standard under homoskedastic errors,
    HC1-robust under heteroskedastic errors, HAC for the dynamic unbalanced
    design (DGP 4).
  \item \textbf{Size}: empirical rejection rate of the two-sided $t$-test of
    $H_0\!:\beta=\beta_0$ at the 5\% nominal level using the raw estimator
    with the standard SE,
    $B^{-1}\sum_{b}\mathbf{1}\bigl\{|\hat\beta_b-\beta_0|/
    \widehat{\mathrm{se}}_b > 1.96\bigr\}$.
    Size distortions mirror the bias of the raw estimator
    ($\mathrm{bias}/\mathrm{SE}\sim\sqrt{N/T}$ in the dynamic designs).
  \item \textbf{$P(\hat r = r)$}: percentage of replications in which the
    information criterion (BIC-type, searched over $r=0,\dots,5$) selects the
    true number of factors $r=2$.
  \item \textbf{Conv.}: fraction of replications in which the alternating
    estimation loop converged within the iteration limit.
\end{itemize}

\noindent Table 5 reports the factor-selection distribution:
$\bar r$ is the mean selected factor number, and
$P(\hat r=2)$, $P(\hat r<2)$, $P(\hat r>2)$ are the percentages of
replications with correct selection, underselection, and overselection.

\clearpage
\section{Homoskedastic errors}

% latex table generated in R 4.5.2 by xtable 1.8-4 package
% Fri Jun 12 12:36:49 2026
\begin{table}[ht]
\centering
\caption{Monte Carlo results: static balanced panel (DGP 1). \emph{N}: units; \emph{T}: periods; OLS/IFE/BC: bias of each estimator; SE/SD: ratio of mean SE to Monte Carlo SD; Cov: empirical 95\% coverage; Size: rejection rate under $H_0$.} 
\label{tab:mc_static_bal}
\begin{tabular}{rrllllllllllllll}
  \hline
 & & \multicolumn{6}{c}{Bias \& RMSE} & & \multicolumn{4}{c}{95\% Coverage} & & & \\
\cline{3-8}\cline{10-13}
 \hline
$N$ & $T$ & OLS & IFE & SD & RMSE & BC & BC.RMSE & SE/SD & Std & Rob & Cl & BC.Std & Size & $P(\hat r=r)$ & Conv. \\ 
  \hline
30 & 10 & 0.394 & 0.021 & 0.076 & 0.079 & 0.021 & 0.079 & 0.773 & .872 & .814 & .843 & .865 & .133 & 0.1 & 1.000 \\ 
  30 & 50 & 0.396 & 0.004 & 0.029 & 0.029 & 0.004 & 0.029 & 0.846 & .896 & .880 & .870 & .898 & .104 & 43.4 & 1.000 \\ 
  30 & 100 & 0.397 & 0.002 & 0.020 & 0.020 & 0.002 & 0.020 & 0.855 & .910 & .888 & .873 & .910 & .091 & 75.8 & 1.000 \\ 
  30 & 200 & 0.400 & 0.001 & 0.014 & 0.015 & 0.001 & 0.015 & 0.840 & .907 & .894 & .885 & .908 & .093 & 83.9 & 1.000 \\ 
  50 & 10 & 0.394 & 0.010 & 0.056 & 0.057 & 0.010 & 0.057 & 0.845 & .900 & .857 & .878 & .900 & .100 & 0.1 & 1.000 \\ 
  50 & 50 & 0.398 & 0.002 & 0.022 & 0.022 & 0.002 & 0.022 & 0.909 & .920 & .912 & .916 & .919 & .080 & 93.8 & 1.000 \\ 
  50 & 100 & 0.400 & 0.001 & 0.015 & 0.015 & 0.001 & 0.015 & 0.906 & .927 & .918 & .915 & .925 & .073 & 100.0 & 1.000 \\ 
  50 & 200 & 0.399 & 0.001 & 0.010 & 0.010 & 0.001 & 0.010 & 0.965 & .946 & .936 & .926 & .944 & .054 & 100.0 & 1.000 \\ 
  100 & 10 & 0.392 & 0.004 & 0.037 & 0.037 & 0.004 & 0.037 & 0.934 & .934 & .903 & .925 & .935 & .066 & 0.0 & 1.000 \\ 
  100 & 50 & 0.398 & 0.001 & 0.015 & 0.015 & 0.001 & 0.015 & 0.970 & .945 & .938 & .939 & .946 & .055 & 99.9 & 1.000 \\ 
  100 & 100 & 0.399 & 0.001 & 0.010 & 0.010 & 0.001 & 0.011 & 0.944 & .937 & .930 & .924 & .935 & .063 & 100.0 & 1.000 \\ 
  100 & 200 & 0.398 & 0.000 & 0.007 & 0.007 & 0.000 & 0.007 & 0.964 & .934 & .928 & .929 & .934 & .066 & 100.0 & 1.000 \\ 
  200 & 10 & 0.391 & 0.002 & 0.026 & 0.026 & 0.002 & 0.026 & 0.959 & .942 & .911 & .940 & .943 & .058 & 0.0 & 1.000 \\ 
  200 & 50 & 0.397 & 0.000 & 0.010 & 0.010 & 0.000 & 0.010 & 0.998 & .954 & .944 & .949 & .955 & .046 & 100.0 & 1.000 \\ 
  200 & 100 & 0.399 & -0.000 & 0.007 & 0.007 & -0.000 & 0.007 & 0.978 & .949 & .942 & .942 & .947 & .051 & 100.0 & 1.000 \\ 
  200 & 200 & 0.401 & 0.000 & 0.005 & 0.005 & 0.000 & 0.005 & 0.950 & .936 & .935 & .933 & .937 & .064 & 100.0 & 1.000 \\ 
   \hline
\end{tabular}
\end{table}

% latex table generated in R 4.5.2 by xtable 1.8-4 package
% Fri Jun 12 12:36:49 2026
\begin{table}[ht]
\centering
\caption{Monte Carlo results: static unbalanced panel (DGP 2, fill = 80\%).} 
\label{tab:mc_static_unb_f80}
\begin{tabular}{rrllllllllllllll}
  \hline
 & & \multicolumn{6}{c}{Bias \& RMSE} & & \multicolumn{4}{c}{95\% Coverage} & & & \\
\cline{3-8}\cline{10-13}
 \hline
$N$ & $T$ & OLS & IFE & SD & RMSE & BC & BC.RMSE & SE/SD & Std & Rob & Cl & BC.Std & Size & $P(\hat r=r)$ & Conv. \\ 
  \hline
30 & 10 & 0.385 & 0.042 & 0.108 & 0.116 & 0.043 & 0.117 & 0.644 & .809 & .740 & .780 & .802 & .197 & 61.7 & 0.974 \\ 
  30 & 50 & 0.397 & 0.005 & 0.032 & 0.033 & 0.005 & 0.033 & 0.916 & .923 & .901 & .887 & .926 & .078 & 85.7 & 1.000 \\ 
  30 & 100 & 0.398 & 0.002 & 0.023 & 0.023 & 0.002 & 0.023 & 0.900 & .913 & .898 & .892 & .913 & .087 & 92.0 & 1.000 \\ 
  30 & 200 & 0.396 & 0.001 & 0.016 & 0.016 & 0.001 & 0.016 & 0.895 & .923 & .903 & .903 & .914 & .077 & 93.1 & 1.000 \\ 
  50 & 10 & 0.392 & 0.018 & 0.070 & 0.072 & 0.019 & 0.073 & 0.806 & .878 & .825 & .867 & .878 & .123 & 75.5 & 0.991 \\ 
  50 & 50 & 0.396 & 0.001 & 0.025 & 0.025 & 0.001 & 0.025 & 0.922 & .924 & .911 & .918 & .923 & .076 & 96.1 & 1.000 \\ 
  50 & 100 & 0.398 & 0.001 & 0.017 & 0.017 & 0.001 & 0.017 & 0.943 & .937 & .928 & .929 & .938 & .063 & 98.7 & 1.000 \\ 
  50 & 200 & 0.398 & 0.000 & 0.012 & 0.012 & 0.000 & 0.012 & 0.917 & .929 & .913 & .914 & .926 & .071 & 99.3 & 1.000 \\ 
  100 & 10 & 0.394 & 0.005 & 0.043 & 0.043 & 0.005 & 0.043 & 0.945 & .936 & .898 & .929 & .934 & .065 & 81.1 & 1.000 \\ 
  100 & 50 & 0.396 & 0.001 & 0.017 & 0.017 & 0.001 & 0.017 & 0.965 & .939 & .926 & .930 & .938 & .061 & 98.5 & 1.000 \\ 
  100 & 100 & 0.397 & -0.000 & 0.012 & 0.012 & -0.000 & 0.012 & 0.942 & .935 & .927 & .931 & .937 & .065 & 99.9 & 1.000 \\ 
  100 & 200 & 0.398 & 0.000 & 0.008 & 0.008 & 0.000 & 0.008 & 0.985 & .941 & .938 & .936 & .943 & .059 & 99.9 & 1.000 \\ 
  200 & 10 & 0.394 & 0.003 & 0.030 & 0.030 & 0.003 & 0.030 & 0.965 & .938 & .905 & .940 & .937 & .062 & 85.6 & 1.000 \\ 
  200 & 50 & 0.399 & 0.000 & 0.012 & 0.012 & 0.000 & 0.012 & 0.968 & .947 & .939 & .945 & .946 & .053 & 99.2 & 1.000 \\ 
  200 & 100 & 0.399 & -0.000 & 0.009 & 0.009 & -0.000 & 0.009 & 0.928 & .923 & .919 & .919 & .922 & .077 & 100.0 & 1.000 \\ 
  200 & 200 & 0.399 & -0.000 & 0.006 & 0.006 & -0.000 & 0.006 & 0.955 & .937 & .935 & .938 & .937 & .063 & 100.0 & 1.000 \\ 
   \hline
\end{tabular}
\end{table}

% latex table generated in R 4.5.2 by xtable 1.8-4 package
% Fri Jun 12 12:36:49 2026
\begin{table}[ht]
\centering
\caption{Monte Carlo results: static unbalanced panel (DGP 2, fill = 60\%).} 
\label{tab:mc_static_unb_f60}
\begin{tabular}{rrllllllllllllll}
  \hline
 & & \multicolumn{6}{c}{Bias \& RMSE} & & \multicolumn{4}{c}{95\% Coverage} & & & \\
\cline{3-8}\cline{10-13}
 \hline
$N$ & $T$ & OLS & IFE & SD & RMSE & BC & BC.RMSE & SE/SD & Std & Rob & Cl & BC.Std & Size & $P(\hat r=r)$ & Conv. \\ 
  \hline
30 & 10 & 0.384 & 0.130 & 0.157 & 0.204 & 0.131 & 0.205 & 0.475 & .536 & .461 & .512 & .529 & .471 & 57.8 & 0.813 \\ 
  30 & 50 & 0.396 & 0.008 & 0.041 & 0.042 & 0.008 & 0.042 & 0.826 & .889 & .853 & .854 & .885 & .111 & 76.9 & 1.000 \\ 
  30 & 100 & 0.398 & 0.005 & 0.028 & 0.028 & 0.005 & 0.028 & 0.854 & .902 & .876 & .871 & .900 & .098 & 89.3 & 1.000 \\ 
  30 & 200 & 0.396 & 0.002 & 0.020 & 0.020 & 0.002 & 0.020 & 0.846 & .896 & .885 & .869 & .895 & .104 & 91.1 & 1.000 \\ 
  50 & 10 & 0.393 & 0.063 & 0.113 & 0.130 & 0.065 & 0.131 & 0.569 & .703 & .627 & .690 & .699 & .301 & 67.6 & 0.863 \\ 
  50 & 50 & 0.397 & 0.002 & 0.030 & 0.030 & 0.002 & 0.030 & 0.896 & .917 & .895 & .892 & .913 & .085 & 90.0 & 1.000 \\ 
  50 & 100 & 0.398 & 0.002 & 0.020 & 0.020 & 0.002 & 0.020 & 0.916 & .935 & .913 & .908 & .932 & .065 & 95.9 & 1.000 \\ 
  50 & 200 & 0.398 & 0.001 & 0.014 & 0.014 & 0.001 & 0.014 & 0.915 & .930 & .916 & .912 & .928 & .070 & 97.8 & 1.000 \\ 
  100 & 10 & 0.395 & 0.026 & 0.059 & 0.065 & 0.027 & 0.065 & 0.813 & .876 & .805 & .862 & .865 & .131 & 70.2 & 0.922 \\ 
  100 & 50 & 0.396 & 0.002 & 0.019 & 0.020 & 0.002 & 0.020 & 0.968 & .937 & .928 & .929 & .937 & .063 & 96.4 & 1.000 \\ 
  100 & 100 & 0.396 & 0.000 & 0.014 & 0.014 & 0.000 & 0.014 & 0.919 & .928 & .920 & .924 & .928 & .072 & 99.8 & 1.000 \\ 
  100 & 200 & 0.398 & 0.000 & 0.009 & 0.009 & 0.000 & 0.009 & 0.974 & .947 & .940 & .934 & .944 & .053 & 100.0 & 1.000 \\ 
  200 & 10 & 0.394 & 0.010 & 0.036 & 0.037 & 0.010 & 0.037 & 0.962 & .949 & .861 & .924 & .937 & .051 & 79.7 & 0.987 \\ 
  200 & 50 & 0.399 & 0.001 & 0.014 & 0.014 & 0.001 & 0.014 & 0.966 & .944 & .928 & .936 & .944 & .056 & 99.1 & 1.000 \\ 
  200 & 100 & 0.399 & -0.000 & 0.010 & 0.010 & -0.000 & 0.010 & 0.939 & .936 & .932 & .930 & .936 & .064 & 99.9 & 1.000 \\ 
  200 & 200 & 0.398 & 0.000 & 0.007 & 0.007 & 0.000 & 0.007 & 0.943 & .940 & .935 & .935 & .938 & .060 & 100.0 & 1.000 \\ 
   \hline
\end{tabular}
\end{table}

% latex table generated in R 4.5.2 by xtable 1.8-4 package
% Fri Jun 12 12:36:49 2026
\begin{table}[ht]
\centering
\caption{Monte Carlo results: dynamic balanced panel (DGP 3). Regressors: $x_{it} = y_{i,t-1}$; bias correction via Moon \& Weidner (2017). $T \geq 50$ only (burn-in requirement). $^*$Cluster-robust SE used as within-unit correlation correction (\texttt{ife()} has no explicit HAC option).} 
\label{tab:mc_dyn_bal}
\begin{tabular}{rrllllllllllllll}
  \hline
 & & \multicolumn{6}{c}{Bias \& RMSE} & & \multicolumn{4}{c}{95\% Coverage} & & & \\
\cline{3-8}\cline{10-13}
 \hline
$N$ & $T$ & OLS & IFE & SD & RMSE & BC & BC.RMSE & SE/SD & Std & Rob & Cl$^*$ & BC.Std & Size & $P(\hat r=r)$ & Conv. \\ 
  \hline
30 & 50 & 0.220 & -0.017 & 0.029 & 0.033 & -0.002 & 0.027 & 0.859 & .854 & .835 & .837 & .936 & .147 & 82.9 & 1.000 \\ 
  30 & 100 & 0.227 & -0.009 & 0.018 & 0.020 & -0.001 & 0.018 & 0.922 & .892 & .880 & .875 & .935 & .108 & 95.4 & 1.000 \\ 
  30 & 200 & 0.228 & -0.003 & 0.013 & 0.013 & 0.000 & 0.013 & 0.928 & .923 & .913 & .895 & .932 & .077 & 98.0 & 1.000 \\ 
  50 & 50 & 0.220 & -0.020 & 0.021 & 0.029 & -0.005 & 0.021 & 0.890 & .802 & .771 & .799 & .927 & .199 & 98.1 & 1.000 \\ 
  50 & 100 & 0.225 & -0.010 & 0.014 & 0.017 & -0.002 & 0.013 & 0.952 & .872 & .859 & .861 & .945 & .128 & 99.9 & 1.000 \\ 
  50 & 200 & 0.227 & -0.005 & 0.009 & 0.011 & -0.001 & 0.009 & 0.964 & .918 & .899 & .903 & .942 & .082 & 100.0 & 1.000 \\ 
  100 & 50 & 0.220 & -0.020 & 0.016 & 0.026 & -0.005 & 0.016 & 0.823 & .653 & .641 & .661 & .914 & .347 & 99.8 & 1.000 \\ 
  100 & 100 & 0.224 & -0.010 & 0.010 & 0.014 & -0.002 & 0.010 & 0.918 & .800 & .792 & .796 & .929 & .200 & 100.0 & 1.000 \\ 
  100 & 200 & 0.228 & -0.005 & 0.006 & 0.008 & -0.001 & 0.006 & 0.975 & .888 & .881 & .882 & .943 & .112 & 100.0 & 1.000 \\ 
  200 & 50 & 0.221 & -0.021 & 0.012 & 0.024 & -0.006 & 0.012 & 0.799 & .424 & .413 & .424 & .885 & .576 & 100.0 & 1.000 \\ 
  200 & 100 & 0.227 & -0.010 & 0.007 & 0.012 & -0.002 & 0.007 & 0.927 & .656 & .646 & .650 & .926 & .344 & 100.0 & 1.000 \\ 
  200 & 200 & 0.228 & -0.005 & 0.005 & 0.007 & -0.001 & 0.005 & 0.935 & .787 & .785 & .784 & .919 & .213 & 100.0 & 1.000 \\ 
   \hline
\end{tabular}
\end{table}

% latex table generated in R 4.5.2 by xtable 1.8-4 package
% Fri Jun 12 12:36:49 2026
\begin{table}[ht]
\centering
\caption{Monte Carlo results: dynamic unbalanced panel (DGP 4, fill = 80\%).} 
\label{tab:mc_dyn_unb_f80}
\begin{tabular}{rrllllllllllllll}
  \hline
 & & \multicolumn{6}{c}{Bias \& RMSE} & & \multicolumn{4}{c}{95\% Coverage} & & & \\
\cline{3-8}\cline{10-13}
 \hline
$N$ & $T$ & OLS & IFE & SD & RMSE & BC & BC.RMSE & SE/SD & Std & Rob & Cov.HAC & BC.HAC & Size & $P(\hat r=r)$ & Conv. \\ 
  \hline
30 & 50 & 0.222 & -0.020 & 0.032 & 0.038 & -0.009 & 0.032 & 0.874 & .860 & .834 & .833 & .897 & .140 & 90.1 & 1.000 \\ 
  30 & 100 & 0.224 & -0.011 & 0.020 & 0.022 & -0.004 & 0.020 & 0.968 & .913 & .890 & .886 & .929 & .087 & 97.3 & 1.000 \\ 
  30 & 200 & 0.227 & -0.005 & 0.013 & 0.014 & -0.002 & 0.013 & 1.003 & .938 & .926 & .927 & .945 & .062 & 98.5 & 1.000 \\ 
  50 & 50 & 0.223 & -0.021 & 0.024 & 0.032 & -0.009 & 0.024 & 0.902 & .812 & .790 & .791 & .898 & .188 & 97.7 & 1.000 \\ 
  50 & 100 & 0.224 & -0.010 & 0.015 & 0.018 & -0.004 & 0.015 & 0.966 & .885 & .871 & .868 & .927 & .115 & 99.7 & 1.000 \\ 
  50 & 200 & 0.228 & -0.005 & 0.011 & 0.012 & -0.001 & 0.011 & 0.950 & .911 & .892 & .895 & .931 & .089 & 100.0 & 1.000 \\ 
  100 & 50 & 0.221 & -0.022 & 0.018 & 0.028 & -0.010 & 0.019 & 0.834 & .665 & .647 & .645 & .855 & .335 & 99.5 & 1.000 \\ 
  100 & 100 & 0.227 & -0.010 & 0.011 & 0.015 & -0.003 & 0.011 & 0.934 & .818 & .806 & .810 & .928 & .182 & 100.0 & 1.000 \\ 
  100 & 200 & 0.229 & -0.005 & 0.008 & 0.009 & -0.002 & 0.008 & 0.941 & .869 & .863 & .863 & .922 & .131 & 100.0 & 1.000 \\ 
  200 & 50 & 0.224 & -0.022 & 0.013 & 0.026 & -0.010 & 0.015 & 0.802 & .451 & .434 & .438 & .812 & .549 & 99.7 & 1.000 \\ 
  200 & 100 & 0.225 & -0.010 & 0.008 & 0.013 & -0.003 & 0.008 & 0.942 & .707 & .693 & .692 & .919 & .293 & 100.0 & 1.000 \\ 
  200 & 200 & 0.227 & -0.005 & 0.005 & 0.007 & -0.001 & 0.005 & 0.972 & .843 & .834 & .836 & .945 & .157 & 100.0 & 1.000 \\ 
   \hline
\end{tabular}
\end{table}

% latex table generated in R 4.5.2 by xtable 1.8-4 package
% Fri Jun 12 12:36:49 2026
\begin{table}[ht]
\centering
\caption{Monte Carlo results: dynamic unbalanced panel (DGP 4, fill = 60\%).} 
\label{tab:mc_dyn_unb_f60}
\begin{tabular}{rrllllllllllllll}
  \hline
 & & \multicolumn{6}{c}{Bias \& RMSE} & & \multicolumn{4}{c}{95\% Coverage} & & & \\
\cline{3-8}\cline{10-13}
 \hline
$N$ & $T$ & OLS & IFE & SD & RMSE & BC & BC.RMSE & SE/SD & Std & Rob & Cov.HAC & BC.HAC & Size & $P(\hat r=r)$ & Conv. \\ 
  \hline
30 & 50 & 0.223 & -0.019 & 0.038 & 0.043 & -0.009 & 0.038 & 0.845 & .875 & .835 & .835 & .881 & .126 & 87.7 & 1.000 \\ 
  30 & 100 & 0.224 & -0.010 & 0.023 & 0.025 & -0.004 & 0.023 & 0.966 & .917 & .893 & .892 & .914 & .084 & 97.2 & 1.000 \\ 
  30 & 200 & 0.226 & -0.005 & 0.016 & 0.017 & -0.002 & 0.016 & 0.985 & .937 & .917 & .918 & .934 & .063 & 99.2 & 1.000 \\ 
  50 & 50 & 0.222 & -0.021 & 0.029 & 0.035 & -0.010 & 0.029 & 0.857 & .836 & .803 & .804 & .876 & .165 & 95.0 & 1.000 \\ 
  50 & 100 & 0.224 & -0.010 & 0.018 & 0.020 & -0.004 & 0.018 & 0.949 & .897 & .883 & .883 & .921 & .103 & 99.5 & 1.000 \\ 
  50 & 200 & 0.228 & -0.005 & 0.012 & 0.013 & -0.001 & 0.012 & 0.967 & .926 & .915 & .915 & .940 & .074 & 100.0 & 1.000 \\ 
  100 & 50 & 0.221 & -0.021 & 0.021 & 0.030 & -0.010 & 0.022 & 0.834 & .724 & .701 & .703 & .872 & .276 & 98.7 & 1.000 \\ 
  100 & 100 & 0.227 & -0.010 & 0.013 & 0.016 & -0.003 & 0.013 & 0.932 & .853 & .841 & .842 & .921 & .147 & 100.0 & 1.000 \\ 
  100 & 200 & 0.228 & -0.005 & 0.009 & 0.010 & -0.002 & 0.009 & 0.946 & .889 & .880 & .878 & .926 & .111 & 100.0 & 1.000 \\ 
  200 & 50 & 0.224 & -0.022 & 0.015 & 0.026 & -0.010 & 0.017 & 0.812 & .557 & .534 & .536 & .805 & .443 & 100.0 & 1.000 \\ 
  200 & 100 & 0.225 & -0.010 & 0.009 & 0.013 & -0.003 & 0.009 & 0.935 & .764 & .746 & .742 & .913 & .236 & 100.0 & 1.000 \\ 
  200 & 200 & 0.227 & -0.005 & 0.006 & 0.008 & -0.001 & 0.006 & 0.955 & .853 & .845 & .843 & .938 & .147 & 100.0 & 1.000 \\ 
   \hline
\end{tabular}
\end{table}

% latex table generated in R 4.5.2 by xtable 1.8-4 package
% Fri Jun 12 12:36:49 2026
\begin{table}[ht]
\centering
\caption{Factor selection accuracy across all DGPs ($r_{\rm true} = 2$, $r_{\max} = 5$).} 
\label{tab:mc_factor_sel}
\begin{tabular}{rlrrllll}
  \hline
DGP & Fill & $N$ & $T$ & $\bar r$ & $P(\hat r=2)$\% & $P(\hat r<2)$\% & $P(\hat r>2)$\% \\ 
  \hline
1 & 80\% & 30 & 10 & 1.00 & 0.1 & 99.9 & 0.0 \\ 
  1 & 80\% & 30 & 50 & 1.43 & 43.4 & 56.6 & 0.0 \\ 
  1 & 80\% & 30 & 100 & 1.76 & 75.8 & 24.2 & 0.0 \\ 
  1 & 80\% & 30 & 200 & 1.84 & 83.9 & 16.1 & 0.0 \\ 
  1 & 80\% & 50 & 10 & 1.00 & 0.1 & 99.9 & 0.0 \\ 
  1 & 80\% & 50 & 50 & 1.94 & 93.8 & 6.2 & 0.0 \\ 
  1 & 80\% & 50 & 100 & 2.00 & 100.0 & 0.0 & 0.0 \\ 
  1 & 80\% & 50 & 200 & 2.00 & 100.0 & 0.0 & 0.0 \\ 
  1 & 80\% & 100 & 10 & 1.00 & 0.0 & 100.0 & 0.0 \\ 
  1 & 80\% & 100 & 50 & 2.00 & 99.9 & 0.1 & 0.0 \\ 
  1 & 80\% & 100 & 100 & 2.00 & 100.0 & 0.0 & 0.0 \\ 
  1 & 80\% & 100 & 200 & 2.00 & 100.0 & 0.0 & 0.0 \\ 
  1 & 80\% & 200 & 10 & 1.00 & 0.0 & 100.0 & 0.0 \\ 
  1 & 80\% & 200 & 50 & 2.00 & 100.0 & 0.0 & 0.0 \\ 
  1 & 80\% & 200 & 100 & 2.00 & 100.0 & 0.0 & 0.0 \\ 
  1 & 80\% & 200 & 200 & 2.00 & 100.0 & 0.0 & 0.0 \\ 
  2 & 60\% & 30 & 10 & 1.69 & 57.8 & 36.5 & 5.7 \\ 
  2 & 60\% & 30 & 50 & 1.77 & 76.9 & 23.1 & 0.0 \\ 
  2 & 60\% & 30 & 100 & 1.89 & 89.3 & 10.7 & 0.0 \\ 
  2 & 60\% & 30 & 200 & 1.91 & 91.1 & 8.9 & 0.0 \\ 
  2 & 60\% & 50 & 10 & 1.71 & 67.6 & 30.8 & 1.6 \\ 
  2 & 60\% & 50 & 50 & 1.90 & 90.0 & 10.0 & 0.0 \\ 
  2 & 60\% & 50 & 100 & 1.96 & 95.9 & 4.1 & 0.0 \\ 
  2 & 60\% & 50 & 200 & 1.98 & 97.8 & 2.2 & 0.0 \\ 
  2 & 60\% & 100 & 10 & 1.70 & 70.2 & 29.8 & 0.0 \\ 
  2 & 60\% & 100 & 50 & 1.96 & 96.4 & 3.6 & 0.0 \\ 
  2 & 60\% & 100 & 100 & 2.00 & 99.8 & 0.2 & 0.0 \\ 
  2 & 60\% & 100 & 200 & 2.00 & 100.0 & 0.0 & 0.0 \\ 
  2 & 60\% & 200 & 10 & 1.80 & 79.7 & 20.3 & 0.0 \\ 
  2 & 60\% & 200 & 50 & 1.99 & 99.1 & 0.9 & 0.0 \\ 
  2 & 60\% & 200 & 100 & 2.00 & 99.9 & 0.1 & 0.0 \\ 
  2 & 60\% & 200 & 200 & 2.00 & 100.0 & 0.0 & 0.0 \\ 
  2 & 80\% & 30 & 10 & 1.70 & 61.7 & 34.1 & 4.2 \\ 
  2 & 80\% & 30 & 50 & 1.86 & 85.7 & 14.2 & 0.1 \\ 
  2 & 80\% & 30 & 100 & 1.93 & 92.0 & 7.8 & 0.2 \\ 
  2 & 80\% & 30 & 200 & 1.95 & 93.1 & 6.5 & 0.4 \\ 
  2 & 80\% & 50 & 10 & 1.78 & 75.5 & 23.5 & 1.0 \\ 
  2 & 80\% & 50 & 50 & 1.97 & 96.1 & 3.8 & 0.1 \\ 
  2 & 80\% & 50 & 100 & 1.99 & 98.7 & 1.3 & 0.0 \\ 
  2 & 80\% & 50 & 200 & 2.00 & 99.3 & 0.5 & 0.2 \\ 
  2 & 80\% & 100 & 10 & 1.81 & 81.1 & 18.8 & 0.1 \\ 
  2 & 80\% & 100 & 50 & 1.99 & 98.5 & 1.5 & 0.0 \\ 
  2 & 80\% & 100 & 100 & 2.00 & 99.9 & 0.1 & 0.0 \\ 
  2 & 80\% & 100 & 200 & 2.00 & 99.9 & 0.0 & 0.1 \\ 
  2 & 80\% & 200 & 10 & 1.86 & 85.6 & 14.4 & 0.0 \\ 
  2 & 80\% & 200 & 50 & 1.99 & 99.2 & 0.8 & 0.0 \\ 
  2 & 80\% & 200 & 100 & 2.00 & 100.0 & 0.0 & 0.0 \\ 
  2 & 80\% & 200 & 200 & 2.00 & 100.0 & 0.0 & 0.0 \\ 
  3 & 80\% & 30 & 50 & 1.83 & 82.9 & 17.1 & 0.0 \\ 
  3 & 80\% & 30 & 100 & 1.95 & 95.4 & 4.6 & 0.0 \\ 
  3 & 80\% & 30 & 200 & 1.98 & 98.0 & 2.0 & 0.0 \\ 
  3 & 80\% & 50 & 50 & 1.98 & 98.1 & 1.9 & 0.0 \\ 
  3 & 80\% & 50 & 100 & 2.00 & 99.9 & 0.1 & 0.0 \\ 
  3 & 80\% & 50 & 200 & 2.00 & 100.0 & 0.0 & 0.0 \\ 
  3 & 80\% & 100 & 50 & 2.00 & 99.8 & 0.2 & 0.0 \\ 
  3 & 80\% & 100 & 100 & 2.00 & 100.0 & 0.0 & 0.0 \\ 
  3 & 80\% & 100 & 200 & 2.00 & 100.0 & 0.0 & 0.0 \\ 
  3 & 80\% & 200 & 50 & 2.00 & 100.0 & 0.0 & 0.0 \\ 
  3 & 80\% & 200 & 100 & 2.00 & 100.0 & 0.0 & 0.0 \\ 
  3 & 80\% & 200 & 200 & 2.00 & 100.0 & 0.0 & 0.0 \\ 
  4 & 60\% & 30 & 50 & 1.88 & 87.7 & 12.3 & 0.0 \\ 
  4 & 60\% & 30 & 100 & 1.97 & 97.2 & 2.8 & 0.0 \\ 
  4 & 60\% & 30 & 200 & 1.99 & 99.2 & 0.8 & 0.0 \\ 
  4 & 60\% & 50 & 50 & 1.95 & 95.0 & 5.0 & 0.0 \\ 
  4 & 60\% & 50 & 100 & 2.00 & 99.5 & 0.5 & 0.0 \\ 
  4 & 60\% & 50 & 200 & 2.00 & 100.0 & 0.0 & 0.0 \\ 
  4 & 60\% & 100 & 50 & 1.99 & 98.7 & 1.3 & 0.0 \\ 
  4 & 60\% & 100 & 100 & 2.00 & 100.0 & 0.0 & 0.0 \\ 
  4 & 60\% & 100 & 200 & 2.00 & 100.0 & 0.0 & 0.0 \\ 
  4 & 60\% & 200 & 50 & 2.00 & 100.0 & 0.0 & 0.0 \\ 
  4 & 60\% & 200 & 100 & 2.00 & 100.0 & 0.0 & 0.0 \\ 
  4 & 60\% & 200 & 200 & 2.00 & 100.0 & 0.0 & 0.0 \\ 
  4 & 80\% & 30 & 50 & 1.90 & 90.1 & 9.9 & 0.0 \\ 
  4 & 80\% & 30 & 100 & 1.97 & 97.3 & 2.7 & 0.0 \\ 
  4 & 80\% & 30 & 200 & 1.99 & 98.5 & 1.5 & 0.0 \\ 
  4 & 80\% & 50 & 50 & 1.98 & 97.7 & 2.3 & 0.0 \\ 
  4 & 80\% & 50 & 100 & 2.00 & 99.7 & 0.3 & 0.0 \\ 
  4 & 80\% & 50 & 200 & 2.00 & 100.0 & 0.0 & 0.0 \\ 
  4 & 80\% & 100 & 50 & 2.00 & 99.5 & 0.5 & 0.0 \\ 
  4 & 80\% & 100 & 100 & 2.00 & 100.0 & 0.0 & 0.0 \\ 
  4 & 80\% & 100 & 200 & 2.00 & 100.0 & 0.0 & 0.0 \\ 
  4 & 80\% & 200 & 50 & 2.00 & 99.7 & 0.3 & 0.0 \\ 
  4 & 80\% & 200 & 100 & 2.00 & 100.0 & 0.0 & 0.0 \\ 
  4 & 80\% & 200 & 200 & 2.00 & 100.0 & 0.0 & 0.0 \\ 
   \hline
\end{tabular}
\end{table}


\clearpage
\section{Heteroskedastic errors}

% latex table generated in R 4.5.2 by xtable 1.8-4 package
% Fri Jun 12 12:37:03 2026
\begin{table}[ht]
\centering
\caption{Monte Carlo results: static balanced panel (DGP 1). \emph{N}: units; \emph{T}: periods; OLS/IFE/BC: bias of each estimator; SE/SD: ratio of mean SE to Monte Carlo SD; Cov: empirical 95\% coverage; Size: rejection rate under $H_0$.} 
\label{tab:mc_static_bal}
\begin{tabular}{rrllllllllllllll}
  \hline
 & & \multicolumn{6}{c}{Bias \& RMSE} & & \multicolumn{4}{c}{95\% Coverage} & & & \\
\cline{3-8}\cline{10-13}
 \hline
$N$ & $T$ & OLS & IFE & SD & RMSE & BC & BC.RMSE & SE/SD & Std & Rob & Cl & BC.Rob & Size & $P(\hat r=r)$ & Conv. \\ 
  \hline
30 & 10 & 0.394 & 0.038 & 0.082 & 0.091 & 0.041 & 0.093 & 0.627 & .788 & .737 & .799 & .724 & .213 & 0.6 & 1.000 \\ 
  30 & 50 & 0.396 & 0.006 & 0.030 & 0.030 & 0.007 & 0.031 & 0.769 & .856 & .864 & .898 & .862 & .144 & 65.1 & 1.000 \\ 
  30 & 100 & 0.397 & 0.004 & 0.021 & 0.021 & 0.004 & 0.021 & 0.771 & .851 & .865 & .921 & .868 & .149 & 87.7 & 1.000 \\ 
  30 & 200 & 0.400 & 0.003 & 0.016 & 0.016 & 0.003 & 0.016 & 0.734 & .844 & .863 & .936 & .871 & .156 & 92.6 & 1.000 \\ 
  50 & 10 & 0.395 & 0.027 & 0.060 & 0.066 & 0.029 & 0.068 & 0.700 & .814 & .778 & .825 & .763 & .187 & 0.8 & 1.000 \\ 
  50 & 50 & 0.398 & 0.003 & 0.022 & 0.022 & 0.003 & 0.022 & 0.860 & .907 & .915 & .933 & .907 & .093 & 97.6 & 1.000 \\ 
  50 & 100 & 0.400 & 0.002 & 0.016 & 0.016 & 0.002 & 0.016 & 0.842 & .897 & .903 & .929 & .917 & .103 & 100.0 & 1.000 \\ 
  50 & 200 & 0.399 & 0.002 & 0.011 & 0.011 & 0.001 & 0.010 & 0.869 & .901 & .920 & .960 & .930 & .099 & 100.0 & 1.000 \\ 
  100 & 10 & 0.392 & 0.016 & 0.040 & 0.044 & 0.018 & 0.045 & 0.771 & .866 & .830 & .870 & .822 & .134 & 0.4 & 1.000 \\ 
  100 & 50 & 0.398 & 0.002 & 0.014 & 0.014 & 0.002 & 0.014 & 0.959 & .938 & .937 & .948 & .933 & .062 & 100.0 & 1.000 \\ 
  100 & 100 & 0.399 & 0.001 & 0.010 & 0.010 & 0.001 & 0.010 & 0.929 & .925 & .929 & .943 & .933 & .075 & 100.0 & 1.000 \\ 
  100 & 200 & 0.398 & 0.000 & 0.007 & 0.007 & 0.000 & 0.007 & 0.939 & .937 & .940 & .960 & .940 & .063 & 100.0 & 1.000 \\ 
  200 & 10 & 0.392 & 0.013 & 0.027 & 0.030 & 0.014 & 0.031 & 0.831 & .877 & .838 & .884 & .826 & .123 & 0.0 & 1.000 \\ 
  200 & 50 & 0.397 & 0.001 & 0.010 & 0.010 & 0.001 & 0.010 & 0.989 & .947 & .943 & .950 & .948 & .053 & 100.0 & 1.000 \\ 
  200 & 100 & 0.400 & 0.001 & 0.007 & 0.007 & 0.000 & 0.007 & 0.996 & .945 & .945 & .953 & .942 & .055 & 100.0 & 1.000 \\ 
  200 & 200 & 0.401 & 0.000 & 0.005 & 0.005 & 0.000 & 0.005 & 0.984 & .945 & .952 & .953 & .950 & .055 & 100.0 & 1.000 \\ 
   \hline
\end{tabular}
\end{table}

% latex table generated in R 4.5.2 by xtable 1.8-4 package
% Fri Jun 12 12:37:03 2026
\begin{table}[ht]
\centering
\caption{Monte Carlo results: static unbalanced panel (DGP 2, fill = 80\%).} 
\label{tab:mc_static_unb_f80}
\begin{tabular}{rrllllllllllllll}
  \hline
 & & \multicolumn{6}{c}{Bias \& RMSE} & & \multicolumn{4}{c}{95\% Coverage} & & & \\
\cline{3-8}\cline{10-13}
 \hline
$N$ & $T$ & OLS & IFE & SD & RMSE & BC & BC.RMSE & SE/SD & Std & Rob & Cl & BC.Rob & Size & $P(\hat r=r)$ & Conv. \\ 
  \hline
30 & 10 & 0.385 & 0.061 & 0.102 & 0.119 & 0.066 & 0.122 & 0.578 & .714 & .654 & .709 & .632 & .286 & 56.2 & 0.957 \\ 
  30 & 50 & 0.396 & 0.009 & 0.034 & 0.035 & 0.010 & 0.035 & 0.817 & .887 & .857 & .879 & .842 & .113 & 83.7 & 1.000 \\ 
  30 & 100 & 0.398 & 0.005 & 0.024 & 0.024 & 0.005 & 0.024 & 0.820 & .885 & .871 & .895 & .866 & .115 & 87.6 & 1.000 \\ 
  30 & 200 & 0.396 & 0.004 & 0.019 & 0.019 & 0.004 & 0.019 & 0.724 & .836 & .823 & .861 & .829 & .164 & 90.2 & 1.000 \\ 
  50 & 10 & 0.393 & 0.042 & 0.078 & 0.089 & 0.046 & 0.091 & 0.624 & .767 & .707 & .765 & .683 & .236 & 69.4 & 0.981 \\ 
  50 & 50 & 0.396 & 0.004 & 0.024 & 0.024 & 0.004 & 0.025 & 0.911 & .920 & .903 & .916 & .899 & .080 & 94.9 & 1.000 \\ 
  50 & 100 & 0.399 & 0.003 & 0.017 & 0.017 & 0.003 & 0.017 & 0.899 & .912 & .902 & .909 & .910 & .088 & 97.2 & 1.000 \\ 
  50 & 200 & 0.398 & 0.001 & 0.012 & 0.013 & 0.001 & 0.012 & 0.869 & .911 & .899 & .917 & .907 & .089 & 98.7 & 1.000 \\ 
  100 & 10 & 0.394 & 0.025 & 0.051 & 0.056 & 0.027 & 0.058 & 0.708 & .818 & .780 & .840 & .762 & .182 & 76.0 & 0.993 \\ 
  100 & 50 & 0.396 & 0.003 & 0.017 & 0.017 & 0.003 & 0.017 & 0.940 & .935 & .925 & .931 & .919 & .065 & 98.3 & 1.000 \\ 
  100 & 100 & 0.397 & 0.001 & 0.012 & 0.012 & 0.001 & 0.012 & 0.923 & .931 & .920 & .929 & .924 & .069 & 99.7 & 1.000 \\ 
  100 & 200 & 0.398 & 0.001 & 0.008 & 0.008 & 0.001 & 0.008 & 0.960 & .945 & .942 & .944 & .948 & .055 & 99.6 & 1.000 \\ 
  200 & 10 & 0.394 & 0.018 & 0.034 & 0.038 & 0.019 & 0.039 & 0.765 & .847 & .790 & .854 & .782 & .153 & 80.9 & 0.998 \\ 
  200 & 50 & 0.399 & 0.002 & 0.012 & 0.012 & 0.002 & 0.012 & 0.965 & .945 & .932 & .943 & .924 & .055 & 99.3 & 1.000 \\ 
  200 & 100 & 0.399 & 0.000 & 0.008 & 0.008 & 0.000 & 0.008 & 0.965 & .941 & .935 & .938 & .935 & .059 & 100.0 & 1.000 \\ 
  200 & 200 & 0.399 & -0.000 & 0.006 & 0.006 & -0.000 & 0.006 & 0.975 & .949 & .948 & .941 & .948 & .051 & 100.0 & 1.000 \\ 
   \hline
\end{tabular}
\end{table}

% latex table generated in R 4.5.2 by xtable 1.8-4 package
% Fri Jun 12 12:37:03 2026
\begin{table}[ht]
\centering
\caption{Monte Carlo results: static unbalanced panel (DGP 2, fill = 60\%).} 
\label{tab:mc_static_unb_f60}
\begin{tabular}{rrllllllllllllll}
  \hline
 & & \multicolumn{6}{c}{Bias \& RMSE} & & \multicolumn{4}{c}{95\% Coverage} & & & \\
\cline{3-8}\cline{10-13}
 \hline
$N$ & $T$ & OLS & IFE & SD & RMSE & BC & BC.RMSE & SE/SD & Std & Rob & Cl & BC.Rob & Size & $P(\hat r=r)$ & Conv. \\ 
  \hline
30 & 10 & 0.384 & 0.110 & 0.128 & 0.169 & 0.114 & 0.172 & 0.481 & .544 & .479 & .534 & .471 & .462 & 54.1 & 0.745 \\ 
  30 & 50 & 0.396 & 0.018 & 0.046 & 0.049 & 0.020 & 0.051 & 0.677 & .826 & .789 & .810 & .776 & .175 & 73.5 & 1.000 \\ 
  30 & 100 & 0.398 & 0.009 & 0.029 & 0.031 & 0.010 & 0.031 & 0.751 & .837 & .810 & .847 & .797 & .163 & 80.8 & 1.000 \\ 
  30 & 200 & 0.396 & 0.006 & 0.023 & 0.024 & 0.006 & 0.024 & 0.666 & .804 & .789 & .844 & .798 & .197 & 80.9 & 1.000 \\ 
  50 & 10 & 0.393 & 0.083 & 0.106 & 0.135 & 0.089 & 0.139 & 0.504 & .602 & .530 & .596 & .510 & .405 & 65.4 & 0.849 \\ 
  50 & 50 & 0.397 & 0.008 & 0.029 & 0.030 & 0.009 & 0.031 & 0.847 & .892 & .868 & .883 & .856 & .108 & 87.6 & 1.000 \\ 
  50 & 100 & 0.398 & 0.005 & 0.021 & 0.022 & 0.005 & 0.022 & 0.837 & .882 & .872 & .891 & .868 & .118 & 93.7 & 1.000 \\ 
  50 & 200 & 0.398 & 0.003 & 0.015 & 0.015 & 0.003 & 0.015 & 0.850 & .895 & .886 & .904 & .883 & .106 & 96.2 & 1.000 \\ 
  100 & 10 & 0.395 & 0.050 & 0.067 & 0.084 & 0.055 & 0.087 & 0.612 & .688 & .642 & .695 & .610 & .312 & 69.5 & 0.904 \\ 
  100 & 50 & 0.396 & 0.005 & 0.019 & 0.020 & 0.006 & 0.020 & 0.930 & .923 & .910 & .921 & .911 & .077 & 94.0 & 1.000 \\ 
  100 & 100 & 0.396 & 0.002 & 0.014 & 0.014 & 0.002 & 0.014 & 0.902 & .923 & .910 & .919 & .910 & .078 & 99.9 & 1.000 \\ 
  100 & 200 & 0.398 & 0.001 & 0.010 & 0.010 & 0.001 & 0.010 & 0.936 & .935 & .930 & .941 & .931 & .065 & 100.0 & 1.000 \\ 
  200 & 10 & 0.394 & 0.032 & 0.042 & 0.053 & 0.035 & 0.055 & 0.717 & .797 & .759 & .835 & .734 & .203 & 77.2 & 0.975 \\ 
  200 & 50 & 0.399 & 0.003 & 0.014 & 0.014 & 0.003 & 0.014 & 0.946 & .931 & .906 & .927 & .909 & .069 & 97.0 & 1.000 \\ 
  200 & 100 & 0.399 & 0.001 & 0.009 & 0.010 & 0.001 & 0.010 & 0.957 & .939 & .933 & .939 & .927 & .061 & 99.9 & 1.000 \\ 
  200 & 200 & 0.399 & 0.000 & 0.007 & 0.007 & 0.000 & 0.007 & 0.953 & .944 & .939 & .941 & .939 & .056 & 100.0 & 1.000 \\ 
   \hline
\end{tabular}
\end{table}

% latex table generated in R 4.5.2 by xtable 1.8-4 package
% Fri Jun 12 12:37:03 2026
\begin{table}[ht]
\centering
\caption{Monte Carlo results: dynamic balanced panel (DGP 3). Regressors: $x_{it} = y_{i,t-1}$; bias correction via Moon \& Weidner (2017). $T \geq 50$ only (burn-in requirement). $^*$Cluster-robust SE used as within-unit correlation correction (\texttt{ife()} has no explicit HAC option).} 
\label{tab:mc_dyn_bal}
\begin{tabular}{rrllllllllllllll}
  \hline
 & & \multicolumn{6}{c}{Bias \& RMSE} & & \multicolumn{4}{c}{95\% Coverage} & & & \\
\cline{3-8}\cline{10-13}
 \hline
$N$ & $T$ & OLS & IFE & SD & RMSE & BC & BC.RMSE & SE/SD & Std & Rob & Cl$^*$ & BC.Rob & Size & $P(\hat r=r)$ & Conv. \\ 
  \hline
30 & 50 & 0.220 & -0.027 & 0.036 & 0.045 & -0.007 & 0.035 & 0.686 & .722 & .759 & .758 & .878 & .278 & 90.6 & 1.000 \\ 
  30 & 100 & 0.227 & -0.014 & 0.023 & 0.027 & -0.004 & 0.023 & 0.726 & .776 & .837 & .814 & .905 & .225 & 98.2 & 1.000 \\ 
  30 & 200 & 0.228 & -0.006 & 0.016 & 0.017 & -0.000 & 0.016 & 0.736 & .826 & .882 & .871 & .924 & .174 & 98.7 & 1.000 \\ 
  50 & 50 & 0.221 & -0.032 & 0.029 & 0.043 & -0.010 & 0.029 & 0.658 & .581 & .664 & .664 & .881 & .419 & 99.1 & 1.000 \\ 
  50 & 100 & 0.225 & -0.015 & 0.018 & 0.023 & -0.004 & 0.018 & 0.726 & .719 & .812 & .792 & .917 & .281 & 100.0 & 1.000 \\ 
  50 & 200 & 0.227 & -0.008 & 0.013 & 0.015 & -0.002 & 0.013 & 0.718 & .767 & .846 & .839 & .907 & .233 & 100.0 & 1.000 \\ 
  100 & 50 & 0.220 & -0.033 & 0.022 & 0.040 & -0.011 & 0.022 & 0.598 & .406 & .498 & .514 & .854 & .594 & 100.0 & 1.000 \\ 
  100 & 100 & 0.224 & -0.017 & 0.013 & 0.021 & -0.005 & 0.014 & 0.683 & .544 & .677 & .674 & .901 & .456 & 100.0 & 1.000 \\ 
  100 & 200 & 0.228 & -0.008 & 0.008 & 0.012 & -0.002 & 0.008 & 0.749 & .691 & .806 & .804 & .936 & .309 & 100.0 & 1.000 \\ 
  200 & 50 & 0.221 & -0.034 & 0.018 & 0.038 & -0.012 & 0.019 & 0.530 & .196 & .255 & .285 & .785 & .804 & 100.0 & 1.000 \\ 
  200 & 100 & 0.227 & -0.017 & 0.010 & 0.019 & -0.005 & 0.010 & 0.648 & .338 & .475 & .476 & .876 & .662 & 100.0 & 1.000 \\ 
  200 & 200 & 0.228 & -0.009 & 0.006 & 0.011 & -0.002 & 0.006 & 0.724 & .521 & .671 & .674 & .919 & .479 & 100.0 & 1.000 \\ 
   \hline
\end{tabular}
\end{table}

% latex table generated in R 4.5.2 by xtable 1.8-4 package
% Fri Jun 12 12:37:03 2026
\begin{table}[ht]
\centering
\caption{Monte Carlo results: dynamic unbalanced panel (DGP 4, fill = 80\%).} 
\label{tab:mc_dyn_unb_f80}
\begin{tabular}{rrllllllllllllll}
  \hline
 & & \multicolumn{6}{c}{Bias \& RMSE} & & \multicolumn{4}{c}{95\% Coverage} & & & \\
\cline{3-8}\cline{10-13}
 \hline
$N$ & $T$ & OLS & IFE & SD & RMSE & BC & BC.RMSE & SE/SD & Std & Rob & Cov.HAC & BC.HAC & Size & $P(\hat r=r)$ & Conv. \\ 
  \hline
30 & 50 & 0.222 & -0.031 & 0.042 & 0.052 & -0.015 & 0.043 & 0.660 & .716 & .745 & .744 & .831 & .285 & 91.6 & 1.000 \\ 
  30 & 100 & 0.224 & -0.015 & 0.025 & 0.029 & -0.005 & 0.025 & 0.754 & .802 & .854 & .852 & .894 & .198 & 97.3 & 1.000 \\ 
  30 & 200 & 0.227 & -0.007 & 0.018 & 0.019 & -0.002 & 0.018 & 0.754 & .821 & .874 & .870 & .906 & .179 & 98.2 & 1.000 \\ 
  50 & 50 & 0.223 & -0.030 & 0.031 & 0.043 & -0.015 & 0.033 & 0.689 & .654 & .724 & .723 & .867 & .346 & 97.7 & 1.000 \\ 
  50 & 100 & 0.224 & -0.015 & 0.019 & 0.025 & -0.005 & 0.020 & 0.747 & .746 & .827 & .831 & .912 & .254 & 99.8 & 1.000 \\ 
  50 & 200 & 0.228 & -0.008 & 0.014 & 0.016 & -0.003 & 0.014 & 0.743 & .796 & .863 & .860 & .907 & .204 & 100.0 & 1.000 \\ 
  100 & 50 & 0.221 & -0.033 & 0.024 & 0.041 & -0.017 & 0.027 & 0.610 & .451 & .546 & .545 & .797 & .550 & 99.7 & 1.000 \\ 
  100 & 100 & 0.227 & -0.016 & 0.014 & 0.022 & -0.006 & 0.015 & 0.713 & .589 & .709 & .703 & .902 & .411 & 100.0 & 1.000 \\ 
  100 & 200 & 0.229 & -0.008 & 0.010 & 0.013 & -0.003 & 0.010 & 0.708 & .691 & .797 & .793 & .910 & .309 & 100.0 & 1.000 \\ 
  200 & 50 & 0.223 & -0.034 & 0.018 & 0.039 & -0.017 & 0.024 & 0.564 & .233 & .315 & .319 & .688 & .767 & 99.8 & 1.000 \\ 
  200 & 100 & 0.225 & -0.016 & 0.011 & 0.019 & -0.006 & 0.011 & 0.674 & .406 & .549 & .546 & .880 & .594 & 100.0 & 1.000 \\ 
  200 & 200 & 0.227 & -0.008 & 0.007 & 0.011 & -0.003 & 0.007 & 0.734 & .565 & .728 & .725 & .933 & .435 & 100.0 & 1.000 \\ 
   \hline
\end{tabular}
\end{table}

% latex table generated in R 4.5.2 by xtable 1.8-4 package
% Fri Jun 12 12:37:03 2026
\begin{table}[ht]
\centering
\caption{Monte Carlo results: dynamic unbalanced panel (DGP 4, fill = 60\%).} 
\label{tab:mc_dyn_unb_f60}
\begin{tabular}{rrllllllllllllll}
  \hline
 & & \multicolumn{6}{c}{Bias \& RMSE} & & \multicolumn{4}{c}{95\% Coverage} & & & \\
\cline{3-8}\cline{10-13}
 \hline
$N$ & $T$ & OLS & IFE & SD & RMSE & BC & BC.RMSE & SE/SD & Std & Rob & Cov.HAC & BC.HAC & Size & $P(\hat r=r)$ & Conv. \\ 
  \hline
30 & 50 & 0.223 & -0.025 & 0.048 & 0.054 & -0.012 & 0.048 & 0.655 & .751 & .748 & .752 & .809 & .250 & 84.3 & 0.999 \\ 
  30 & 100 & 0.224 & -0.012 & 0.030 & 0.032 & -0.004 & 0.030 & 0.740 & .836 & .866 & .865 & .901 & .164 & 92.9 & 1.000 \\ 
  30 & 200 & 0.227 & -0.006 & 0.021 & 0.022 & -0.001 & 0.021 & 0.746 & .840 & .875 & .872 & .884 & .160 & 97.5 & 1.000 \\ 
  50 & 50 & 0.222 & -0.026 & 0.036 & 0.045 & -0.013 & 0.037 & 0.666 & .720 & .755 & .754 & .833 & .281 & 94.2 & 1.000 \\ 
  50 & 100 & 0.224 & -0.014 & 0.023 & 0.026 & -0.005 & 0.023 & 0.739 & .784 & .829 & .825 & .895 & .216 & 99.3 & 1.000 \\ 
  50 & 200 & 0.228 & -0.007 & 0.016 & 0.017 & -0.002 & 0.016 & 0.752 & .819 & .875 & .870 & .909 & .181 & 99.8 & 1.000 \\ 
  100 & 50 & 0.221 & -0.030 & 0.027 & 0.040 & -0.016 & 0.030 & 0.630 & .549 & .629 & .627 & .807 & .451 & 98.9 & 1.000 \\ 
  100 & 100 & 0.227 & -0.015 & 0.016 & 0.022 & -0.006 & 0.017 & 0.710 & .667 & .763 & .765 & .885 & .333 & 100.0 & 1.000 \\ 
  100 & 200 & 0.229 & -0.008 & 0.011 & 0.014 & -0.002 & 0.011 & 0.720 & .759 & .849 & .848 & .914 & .241 & 100.0 & 1.000 \\ 
  200 & 50 & 0.223 & -0.031 & 0.020 & 0.037 & -0.017 & 0.025 & 0.588 & .349 & .439 & .440 & .724 & .651 & 99.5 & 1.000 \\ 
  200 & 100 & 0.225 & -0.016 & 0.012 & 0.020 & -0.006 & 0.013 & 0.676 & .519 & .645 & .646 & .873 & .482 & 100.0 & 1.000 \\ 
  200 & 200 & 0.227 & -0.008 & 0.008 & 0.011 & -0.003 & 0.008 & 0.734 & .640 & .784 & .778 & .924 & .360 & 100.0 & 1.000 \\ 
   \hline
\end{tabular}
\end{table}

% latex table generated in R 4.5.2 by xtable 1.8-4 package
% Fri Jun 12 12:37:03 2026
\begin{table}[ht]
\centering
\caption{Factor selection accuracy across all DGPs ($r_{\rm true} = 2$, $r_{\max} = 5$).} 
\label{tab:mc_factor_sel}
\begin{tabular}{rlrrllll}
  \hline
DGP & Fill & $N$ & $T$ & $\bar r$ & $P(\hat r=2)$\% & $P(\hat r<2)$\% & $P(\hat r>2)$\% \\ 
  \hline
1 & 80\% & 30 & 10 & 1.01 & 0.6 & 99.4 & 0.0 \\ 
  1 & 80\% & 30 & 50 & 1.65 & 65.1 & 34.9 & 0.0 \\ 
  1 & 80\% & 30 & 100 & 1.88 & 87.7 & 12.3 & 0.0 \\ 
  1 & 80\% & 30 & 200 & 1.93 & 92.6 & 7.4 & 0.0 \\ 
  1 & 80\% & 50 & 10 & 1.01 & 0.8 & 99.2 & 0.0 \\ 
  1 & 80\% & 50 & 50 & 1.98 & 97.6 & 2.4 & 0.0 \\ 
  1 & 80\% & 50 & 100 & 2.00 & 100.0 & 0.0 & 0.0 \\ 
  1 & 80\% & 50 & 200 & 2.00 & 100.0 & 0.0 & 0.0 \\ 
  1 & 80\% & 100 & 10 & 1.00 & 0.4 & 99.6 & 0.0 \\ 
  1 & 80\% & 100 & 50 & 2.00 & 100.0 & 0.0 & 0.0 \\ 
  1 & 80\% & 100 & 100 & 2.00 & 100.0 & 0.0 & 0.0 \\ 
  1 & 80\% & 100 & 200 & 2.00 & 100.0 & 0.0 & 0.0 \\ 
  1 & 80\% & 200 & 10 & 1.00 & 0.0 & 100.0 & 0.0 \\ 
  1 & 80\% & 200 & 50 & 2.00 & 100.0 & 0.0 & 0.0 \\ 
  1 & 80\% & 200 & 100 & 2.00 & 100.0 & 0.0 & 0.0 \\ 
  1 & 80\% & 200 & 200 & 2.00 & 100.0 & 0.0 & 0.0 \\ 
  2 & 60\% & 30 & 10 & 1.63 & 54.1 & 41.5 & 4.4 \\ 
  2 & 60\% & 30 & 50 & 1.75 & 73.5 & 26.0 & 0.5 \\ 
  2 & 60\% & 30 & 100 & 1.81 & 80.8 & 19.2 & 0.0 \\ 
  2 & 60\% & 30 & 200 & 1.81 & 80.9 & 19.1 & 0.0 \\ 
  2 & 60\% & 50 & 10 & 1.71 & 65.4 & 31.9 & 2.7 \\ 
  2 & 60\% & 50 & 50 & 1.88 & 87.6 & 12.4 & 0.0 \\ 
  2 & 60\% & 50 & 100 & 1.94 & 93.7 & 6.3 & 0.0 \\ 
  2 & 60\% & 50 & 200 & 1.96 & 96.2 & 3.8 & 0.0 \\ 
  2 & 60\% & 100 & 10 & 1.71 & 69.5 & 29.8 & 0.7 \\ 
  2 & 60\% & 100 & 50 & 1.94 & 94.0 & 6.0 & 0.0 \\ 
  2 & 60\% & 100 & 100 & 2.00 & 99.9 & 0.1 & 0.0 \\ 
  2 & 60\% & 100 & 200 & 2.00 & 100.0 & 0.0 & 0.0 \\ 
  2 & 60\% & 200 & 10 & 1.77 & 77.2 & 22.8 & 0.0 \\ 
  2 & 60\% & 200 & 50 & 1.97 & 97.0 & 3.0 & 0.0 \\ 
  2 & 60\% & 200 & 100 & 2.00 & 99.9 & 0.1 & 0.0 \\ 
  2 & 60\% & 200 & 200 & 2.00 & 100.0 & 0.0 & 0.0 \\ 
  2 & 80\% & 30 & 10 & 1.71 & 56.2 & 37.3 & 6.5 \\ 
  2 & 80\% & 30 & 50 & 1.89 & 83.7 & 14.6 & 1.7 \\ 
  2 & 80\% & 30 & 100 & 1.94 & 87.6 & 10.8 & 1.6 \\ 
  2 & 80\% & 30 & 200 & 1.95 & 90.2 & 8.6 & 1.2 \\ 
  2 & 80\% & 50 & 10 & 1.75 & 69.4 & 27.9 & 2.6 \\ 
  2 & 80\% & 50 & 50 & 1.99 & 94.9 & 4.0 & 1.1 \\ 
  2 & 80\% & 50 & 100 & 2.00 & 97.2 & 2.0 & 0.8 \\ 
  2 & 80\% & 50 & 200 & 2.02 & 98.7 & 0.5 & 0.8 \\ 
  2 & 80\% & 100 & 10 & 1.79 & 76.0 & 22.4 & 1.6 \\ 
  2 & 80\% & 100 & 50 & 1.99 & 98.3 & 1.6 & 0.1 \\ 
  2 & 80\% & 100 & 100 & 2.00 & 99.7 & 0.1 & 0.2 \\ 
  2 & 80\% & 100 & 200 & 2.01 & 99.6 & 0.0 & 0.4 \\ 
  2 & 80\% & 200 & 10 & 1.81 & 80.9 & 18.8 & 0.3 \\ 
  2 & 80\% & 200 & 50 & 2.00 & 99.3 & 0.6 & 0.1 \\ 
  2 & 80\% & 200 & 100 & 2.00 & 100.0 & 0.0 & 0.0 \\ 
  2 & 80\% & 200 & 200 & 2.00 & 100.0 & 0.0 & 0.0 \\ 
  3 & 80\% & 30 & 50 & 1.91 & 90.6 & 9.4 & 0.0 \\ 
  3 & 80\% & 30 & 100 & 1.98 & 98.2 & 1.8 & 0.0 \\ 
  3 & 80\% & 30 & 200 & 1.99 & 98.7 & 1.3 & 0.0 \\ 
  3 & 80\% & 50 & 50 & 1.99 & 99.1 & 0.9 & 0.0 \\ 
  3 & 80\% & 50 & 100 & 2.00 & 100.0 & 0.0 & 0.0 \\ 
  3 & 80\% & 50 & 200 & 2.00 & 100.0 & 0.0 & 0.0 \\ 
  3 & 80\% & 100 & 50 & 2.00 & 100.0 & 0.0 & 0.0 \\ 
  3 & 80\% & 100 & 100 & 2.00 & 100.0 & 0.0 & 0.0 \\ 
  3 & 80\% & 100 & 200 & 2.00 & 100.0 & 0.0 & 0.0 \\ 
  3 & 80\% & 200 & 50 & 2.00 & 100.0 & 0.0 & 0.0 \\ 
  3 & 80\% & 200 & 100 & 2.00 & 100.0 & 0.0 & 0.0 \\ 
  3 & 80\% & 200 & 200 & 2.00 & 100.0 & 0.0 & 0.0 \\ 
  4 & 60\% & 30 & 50 & 1.85 & 84.3 & 15.5 & 0.2 \\ 
  4 & 60\% & 30 & 100 & 1.93 & 92.9 & 7.1 & 0.0 \\ 
  4 & 60\% & 30 & 200 & 1.98 & 97.5 & 2.5 & 0.0 \\ 
  4 & 60\% & 50 & 50 & 1.94 & 94.2 & 5.8 & 0.0 \\ 
  4 & 60\% & 50 & 100 & 1.99 & 99.3 & 0.7 & 0.0 \\ 
  4 & 60\% & 50 & 200 & 2.00 & 99.8 & 0.2 & 0.0 \\ 
  4 & 60\% & 100 & 50 & 1.99 & 98.9 & 1.1 & 0.0 \\ 
  4 & 60\% & 100 & 100 & 2.00 & 100.0 & 0.0 & 0.0 \\ 
  4 & 60\% & 100 & 200 & 2.00 & 100.0 & 0.0 & 0.0 \\ 
  4 & 60\% & 200 & 50 & 2.00 & 99.5 & 0.5 & 0.0 \\ 
  4 & 60\% & 200 & 100 & 2.00 & 100.0 & 0.0 & 0.0 \\ 
  4 & 60\% & 200 & 200 & 2.00 & 100.0 & 0.0 & 0.0 \\ 
  4 & 80\% & 30 & 50 & 1.92 & 91.6 & 8.3 & 0.1 \\ 
  4 & 80\% & 30 & 100 & 1.97 & 97.3 & 2.7 & 0.0 \\ 
  4 & 80\% & 30 & 200 & 1.98 & 98.2 & 1.8 & 0.0 \\ 
  4 & 80\% & 50 & 50 & 1.98 & 97.7 & 2.1 & 0.2 \\ 
  4 & 80\% & 50 & 100 & 2.00 & 99.8 & 0.2 & 0.0 \\ 
  4 & 80\% & 50 & 200 & 2.00 & 100.0 & 0.0 & 0.0 \\ 
  4 & 80\% & 100 & 50 & 2.00 & 99.7 & 0.3 & 0.0 \\ 
  4 & 80\% & 100 & 100 & 2.00 & 100.0 & 0.0 & 0.0 \\ 
  4 & 80\% & 100 & 200 & 2.00 & 100.0 & 0.0 & 0.0 \\ 
  4 & 80\% & 200 & 50 & 2.00 & 99.8 & 0.2 & 0.0 \\ 
  4 & 80\% & 200 & 100 & 2.00 & 100.0 & 0.0 & 0.0 \\ 
  4 & 80\% & 200 & 200 & 2.00 & 100.0 & 0.0 & 0.0 \\ 
   \hline
\end{tabular}
\end{table}


\end{document}
